% Impaginazione del PDF italiano delle Diatribe.
% Il libro contiene solo la traduzione italiana: il greco e l'inglese stanno sul
% sito. Fonti e crediti sono in fondo.

\usepackage{microtype}
\usepackage{enumitem}
\usepackage{setspace}
\usepackage{graphicx}
\usepackage{xcolor}
\usepackage{eso-pic}
\usepackage{fontspec}

% Palatino Linotype ha il greco politonico ed è la stessa faccia usata dal sito.
% Se manca, xelatex ripiega sul font di default.
\IfFontExistsTF{Palatino Linotype}{\setmainfont{Palatino Linotype}}{}

% Titoli nello stesso carattere del testo: in un libro il contrasto con un
% carattere bastone stona.
\addtokomafont{disposition}{\rmfamily}

\definecolor{eptealdark}{HTML}{1F5F5C}
\definecolor{epink}{HTML}{2A2A26}
\definecolor{epcream}{HTML}{F3EBD8}

% I colori dei collegamenti sono impostati da Quarto (linkcolor/urlcolor nel
% front matter di pdf_book_it.qmd) e si riferiscono ai nomi definiti qui sopra.

% Copertina a piena pagina, prima dell'indice: l'illustrazione fa da sfondo e il
% titolo è composto sopra, nella parte alta che l'immagine lascia libera.
% Gli hook con l'asterisco valgono solo per questa pagina.
\newcommand{\epcover}[1]{%
  \thispagestyle{empty}%
  \AddToShipoutPictureBG*{%
    \AtPageLowerLeft{\color{epcream}\rule{\paperwidth}{\paperheight}}%
  }%
  \AddToShipoutPictureFG*{%
    \AtPageCenter{%
      \makebox(0,0){%
        \includegraphics[width=\paperwidth,height=\paperheight,keepaspectratio]{#1}%
      }%
    }%
    \AtPageUpperLeft{%
      \raisebox{-0.15\paperheight}{%
        \makebox[\paperwidth][c]{%
          \begin{minipage}{0.88\paperwidth}
            \centering
            {\color{eptealdark}\fontsize{17}{20}\selectfont\bfseries
             \addfontfeature{LetterSpace=18}EPITTETO\par}
            \vspace{0.55em}
            {\color{epink}\fontsize{48}{51}\selectfont\bfseries Diatribe\par}
            \vspace{0.7em}
            {\color{black!62}\fontsize{16}{19}\selectfont\itshape
             raccolte da Arriano\par}
          \end{minipage}%
        }%
      }%
    }%
  }%
  \null\clearpage%
}

\setstretch{1.04}
\setlength{\parindent}{0pt}
\setlength{\parskip}{0.34em}

% Sezioni numerate: il numero del versetto vive in un margine a sinistra (più
% piccolo del testo, come nel gutter della versione HTML) e tutto il testo della
% sezione — righe a capo e capoversi successivi inclusi — si allinea a un
% rientro fisso.
\newlength{\verseindent}
\setlength{\verseindent}{2.4em}
\newcommand{\verselabel}[1]{{\footnotesize\bfseries #1}}
\newenvironment{versetto}[1]{%
  \par
  \begingroup
  \setlength{\parindent}{0pt}%
  \leftskip=\verseindent
  \noindent
  \llap{\makebox[\verseindent][l]{\verselabel{#1}}}%
  \ignorespaces
}{%
  \par
  \endgroup
}

% L'indice sta subito dopo la copertina; la pagina successiva è riservata alla
% nota introduttiva.
\let\epstandardtableofcontents\tableofcontents
\renewcommand{\tableofcontents}{%
  \begingroup
  \setstretch{0.96}%
  \epstandardtableofcontents
  \endgroup%
  \clearpage%
}

\setlist[description]{
  style=nextline,
  leftmargin=0pt,
  labelsep=0pt,
  itemsep=0.35em,
  parsep=0pt,
  topsep=0.2em
}
